%% LyX 2.0.7.1 created this file.  For more info, see http://www.lyx.org/.
%% Do not edit unless you really know what you are doing.
\documentclass[11pt,english]{article}
\usepackage[latin9]{inputenc}
\usepackage{geometry}
\geometry{verbose,tmargin=1in,bmargin=1in,lmargin=1in,rmargin=1in}
\usepackage{float}
\usepackage{textcomp}
\usepackage{amstext}
\usepackage{graphicx}
\usepackage{esint}

\makeatletter

%%%%%%%%%%%%%%%%%%%%%%%%%%%%%% LyX specific LaTeX commands.
\DeclareRobustCommand{\greektext}{%
  \fontencoding{LGR}\selectfont\def\encodingdefault{LGR}}
\DeclareRobustCommand{\textgreek}[1]{\leavevmode{\greektext #1}}
\DeclareFontEncoding{LGR}{}{}
\DeclareTextSymbol{\~}{LGR}{126}
\newcommand{\lyxmathsym}[1]{\ifmmode\begingroup\def\b@ld{bold}
  \text{\ifx\math@version\b@ld\bfseries\fi#1}\endgroup\else#1\fi}

%% Because html converters don't know tabularnewline
\providecommand{\tabularnewline}{\\}
%% A simple dot to overcome graphicx limitations
\newcommand{\lyxdot}{.}


\makeatother

\usepackage{babel}
\begin{document}

\title{Digital Lab: Sampling, Fourier Transforms, Mixers, and Down-Converters}


\author{By Patrick Kantorski}


\date{March 7, 2014}
\maketitle
\begin{abstract}
This report was written for a radio astronomy undergraduate course
at UC Berkeley. For this lab, my partners and I examined and tested
various concepts involved with digital signal processing. From Nyquist's
criterion to signal recovery, we performed experiments and collected
data. Additionally, digital and analog mixing (along with SSB and
DSB mixing) was examined. We also provided a design for our FIR filter
and test it by using a digital down converter.
\end{abstract}

\section{Introduction}

Analog signals can be converted to digital signals, but only at a
specified sample rate and for a specified period of time. Many experiments
in Radio Astronomy require recording and extracting information from
radio signals. Although this is almost always done digitally, some
parts of the signal mixing and filtering process can be performed
with analog components. Signals possess many important properties
and can even be mixed together and filtered out.

For this lab, my partners and I examined these concepts by using and
implementing digital signal processing with a CPU based computer,
function generators, ZAD-1 signal mixer, and an FPGA based computer
called the ROACH. Plots and experimental results derived required
extensive coding performed through the programming language Python.
From Nyquist's criterion, signal sampling can be seen to have inherit
limitations which must be considered to properly sample data. Additionally,
we explored both analog and digital methods of signal mixing and compared
the advantages and disadvantages of each. Techniques like as Fourier
filtering can be used on signals in order to recover specific frequencies.
Toward the end of this lab, my partners and I examined and programmed
a FIR filter using the ROACH similar to that of a bandpass. 


\section{Materials \& Methods}


\subsection{Nyquist's Criterion}

Digital data collection functions by collecting sets of data points
at a particular sample rate or sample frequency. One of the fundamental
underlying concepts involved in digital data sampling is Nyquist's
criterion. Nyquist's criterion for digital data collection states
that there is a minimum sample frequency required to digitally sample
a waveform in facsimile. To reproduce a sampled waveform, at least
two data points must be taken per period. The Nyquist frequency is
shown in Equation \ref{eq:nyquist}.

\begin{equation}
v_{sample}=2\times v_{signal}\label{eq:nyquist}
\end{equation}
Because sampled points are separated in time, oscillations in the
signal that occur between two data points will not be recorded. If
Nyquist's criterion is not satisfied during data collection, aliasing
will inevitably occur. Aliasing is when a sampled waveform facsimiles
a lower frequency than what is being measured. An example of aliasing
is shown in Figure 1.
\begin{figure}[H]
\begin{centering}
\includegraphics[scale=0.36]{/Users/ppkantorski/Documents/Radio_Astronomy/Lab_2/Pictures/aliasing}
\par\end{centering}

\caption{Although a clear sine wave is produced from sampling the signal shown
above in red, there are several points missed in-between each of the
data points. This misrepresentation of the original signal is called
aliasing. }
\end{figure}


To demonstrate Nyquist's criterion and aliasing, my partners and I
used the Pulsar sampler card, a computer, and a SRS Function Generator
to take digital waveform data samples of sine waves generated at different
signal frequencies. Additionally, we set our sampling rate to $v_{sample}=1MHz$.
Because the Nyquist frequency, following Equation \ref{eq:nyquist},
was precisely$v_{signal}=0.5MHz$, for observational purposes, we
took data at signals ranging from $v_{signal}=0.1MHz$ to $0.9MHz$.
The sampled waveforms for these signals are shown in Figures 2 and
3.


\subsection{Fourier Analysis}

A Fourier transform is a mathematical transformation used to transform
signals between the time or spatial domain and frequency domain. In
signal and data processing, Fourier transforms can be used to generate
a power spectrum of a digital signal. Besides application, Fourier
transforms are also very useful because sine waves are quite easy
to work with mathematically and form a basis over all function spaces.
Fourier transforms also distinguish negative and positive frequencies
in a spectrum. A negative frequency is the same as a positive frequency
except that instead of the cosine leading the sine by $\frac{\pi}{2}$,
it trails by $\frac{\pi}{2}$ (see Equation \ref{eq:complex}). For
functions with real-values, positive and negative frequencies contain
the same information. In complex-valued functions, positive and negative
frequencies contain independent information, making the transformation
not centered about the origin.

\begin{equation}
e^{i\omega t}=\cos(\omega t)+i\sin(\omega t)\label{eq:complex}
\end{equation}


\begin{eqnarray}
E(v)=\frac{1}{T}\int_{-T/2}^{T/2}E(t)e^{2\pi jvt}dt, &  & E(t)=\frac{1}{F}\int_{-F/2}^{F/2}E(v)e^{-2\pi jvt}dt\label{eq:fourier}
\end{eqnarray}


Fourier transforms are also invertible, allowing one to go from either
go from the time to the frequency domain or from the frequency domain
back to the time domain using the inverse transformation. In the proper
analytic formulation (see Equation \ref{eq:fourier}), the boundaries
of integration for the forward and inverse FT are from negative to
positive infinity. However in application, these boundaries cannot
be infinite. 

Two types of Fourier transforms were implemented and examined in this
lab. One of them is known as the Discrete Fourier Transform or DFT.
A DFT converts a finite list of equally spaced data samples taken
from a function or signal into the list of coefficients of a finite
combination of complex sine waves that has those same sample values,
ordered by their frequencies. Since DFTs function with a finite amount
of data, they can be implemented in computers by numerical algorithms
or even dedicated hardware. Using the programming language Python,
a DFT was constructed by my partners and I to view the power spectrum
of the Nyquist sampling data. 

DFTs have N input numbers and N output numbers, making and the number
of calculations proportional to $N^{2}$. However, for larger N values,
the calculation can take a long time. Because of this, DFTs are inherently
limited. To get around this problem, the Fast Fourier Transform, also
known as the FFT, was developed as to perform Fourier transforms a
lot quicker. Unlike a DFT, in a FFT, the number of calculations is
proportional to $N\log_{2}N$. Since this makes larger N value transforms
more practical, FFTs are typically used more often than DFTs. 


\subsection{Signal Mixing}

While working with radio transmissions, desired signals are often
at frequencies higher than what digital sampling tools are capable
of handling. By using a heterodyne mixer, radio signals can be shifted
down to a lower frequency so that they can be properly sampled. A
local oscillator converts an incoming signal to a different frequency
by mixing the signal with a local oscillated (or LO) signal. When
two signals are mixed, their values are multiplied in the time domain.
Additionally, this is equivalent to taking the convolution of the
signals in the frequency domain (see Equation \ref{eq:convolution}). 

\begin{equation}
f(t)*g(t)=\int f(\omega)g(\omega)e^{i\omega t}d\omega\label{eq:convolution}
\end{equation}


The convolution is a mathematical operation used in applications ranging
from photo editing to radio astronomy. In simple terms, taking the
convolution of two signals effectively ``blurs'' or smears one over
the other. In this lab, my partners and I examined both analog and
digital mixing. 

To analog mixing our signals, we used a Mini-Circuits ZAD-1 and chose
a specific LO frequency to mix along with another sine wave set at
a different frequency (see Equation \ref{eq:analog_mix}). The ZAD-1
is a balanced mixer with identical \textquotedblleft{}R\textquotedblright{}
and \textquotedblleft{}L\textquotedblright{} ports. It also does not
couple to DC or very low frequencies. The \textquotedblleft{}I\textquotedblright{}
port is coupled differently and will handle voltages all the way down
to and including DC. The mixing process functions no matter which
two ports are used as inputs. We also had to make sure that we had
enough samples per period to generate a reasonable facsimile of the
sine wave. This forced us to sample above at least twice the Nyquist
frequency and set the number of points sampled to be large enough
to give produce at least a few periods of the ``slowest'' sine wave.

\begin{equation}
v_{signal}=v_{lo}\pm\delta v\label{eq:analog_mix}
\end{equation}


A single sideband mixer, or SSB, requires mixing with pure sine waves
which have both a real and imaginary part. Because of this, SSB mixing
removes half of the sidebands created. A double sideband mixer, or
DSB, splits the signal into a lower sideband and an upper sideband.
We first took our data samples at the mixed signal frequency $v_{lo}+\delta v$.
and $v_{lo}-\delta v$. Since both signals were by the SRS Function
Generator, they were inherently real and consisted of both $e^{i\lyxmathsym{\textgreek{w}}t}$
and $e^{-i\text{\textgreek{w}}t}$ components (see Equation \ref{eq:complex}).
This made the mixed output of our analog mixed signal composed of
four different frequencies; $v_{lo}+\delta v$, $v_{lo}-\delta v$,
$-v_{lo}+\delta v$, and $-v_{lo}-\delta v$. 


\subsection{Using the ROACH}

The Reconfigurable Open Architecture for Computing Hardware (or ROACH)
is an open-source board developed at the University of California,
Berkeley by Dan Werthimer and Aaron Parsons from the CASPER collaboration.
The circuit board of the ROACH used consists of a Virtex-5 FPGA made
by Xilinx along with a CPU. The CPU allowing users to SSH directly
to the board using a special linux kernel designed by CASPER. This
kernel also allows for automatic interface support between programing
the FPGA. 
\begin{figure}[H]
\begin{centering}
\includegraphics[scale=0.55]{/Users/ppkantorski/Documents/Radio_Astronomy/Lab_2/Pictures/roach}
\par\end{centering}

\caption{The picture above shows the data sampler program design used on the
ROACH, written in the programming language Simulink. This design compiles
data into the \textquotedblleft{}adc snaps\textquotedblright{} BOF
files which are used during digital mixing and sampling.}
\end{figure}


The ROACH is system constructed from a computer running a Linux kernel
and a Field Programmable Gate Array (or FPGA) processor. Unlike CPU
processors, FPGAs can be configured to perform multiple tasks simultaneously
by a series of programable logic gates. In this lab, our FPGA processor
was synchronized by a 200MHz clock. By connecting remotely, my partners
and I programed the ROACH with configurations set to take data samples
through altering the FPGA\textquoteright{}s registers. Data taken
from the ADC output is stored in the form of block random access memory
(or BRAM). After being written, data then can be copied from the ROACH
for speculation and analysis.


\subsection{FIR Filtering}

The Finite Impulse Response filter (or FIR) works by taking the convolution
of two signals in the time domain. Fourier filtering works well with
filtering signals after they have been sampled, but not so much with
signals in real-time. However, the digital implementation can still
be achieved through programming a FPGA. By taking the Fourier transform
of a signal, multiplying (or convolving) the FFTs by a filter, and
then inverting transform to get the filtered signal, real-time digital
filtration can be achieved (see Figure 3). 
\begin{figure}[H]
\begin{centering}
.\includegraphics[scale=0.34]{/Users/ppkantorski/Documents/Radio_Astronomy/Lab_2/Pictures/FIR}
\par\end{centering}

\caption{Figure above is a schematic diagram of how a FIR filter works. Coefficients
or ``taps'' are represent by b$_{N}$.}
\end{figure}


Real-time digital filters perform a convolution on samples received
as the signal is sampled. The implementation of an FIR filter on the
FPGA allows us to convolve an incoming signal with other waveforms
in the time domain. Because the filter my partners and I used has
8 taps, we decided to choose 8 real and 8 imaginary coefficients for
our filter. We did this by writing the coefficient registers to the
FPGA. As the signal is sampled, the FPGA multiplies about 12 of the
last eight points from the ADC with their corresponding coefficients
and then sums them. As the clock cycles, the next point sampled is
fed in and the process repeats. The values are then summed up to form
a new ``convolved'' waveform. Since the sums are largest when the
signal falls along the filter, the FIR filters out portions of the
signal that are similar to the filter's shape and zeros out the rest.


\section{Observations \& Data Analysis}


\subsection{Digital Data Sampling}

To demonstrate aliasing and Nyquist's criterion, data samples were
collected from a SRS Function Generator, over the command-line in
Python. The number of samples collected was 256 with a sample frequency
of 1MHz while the voltage was set to precisely $1V{}_{pp}$. 
\begin{figure}[H]
\begin{centering}
\includegraphics[scale=0.34]{/Users/ppkantorski/Documents/Radio_Astronomy/Lab_2/Pictures/1\lyxdot 2\lyxdot 1}
\par\end{centering}

\caption{The graph above shows data samples of 1MHz, $1V_{pp}$ sine wave collected
through the command-line using Python and a local oscillator and a
Pulsar sampler. }
\end{figure}
\begin{figure}[H]
\begin{centering}
\includegraphics[scale=0.34]{/Users/ppkantorski/Documents/Radio_Astronomy/Lab_2/Pictures/1\lyxdot 4\lyxdot 2}
\par\end{centering}

\caption{The graph above shows the DFT of the data samples from Figure 3. }
\end{figure}


As shown in Figure 4, the signal viewed ranged from $v_{signal}=0.1MHz$
to $0.9MHz$. For signals above at the Nyquist frequency $v_{signal}=0.5MHz$,
the sampled waveform does not have any rising or lowering fluctuations
between the high and low points on the graph. Additionally, the voltage
values at the Nyquist frequency are consistently at about the same
value and same for both high and low values. This is because the incoming
signal is sampled at twice per period. Because there is only are only
high and low data points, the shape of the waveform is triangular
as opposed to a clear facsimile of the sinusoid. Nevertheless, the
0.5MHz frequency is still apparent in both its time domain representation
and Fourier transformation (see Figure 5). 

When the signal frequency is above 0.5MHz, the effects of aliasing
becomes distinctive. After falling below the Nyquist frequency, the
data plots begin to revert. At the signal frequency of 0.9MHz, a clear
sine wave can be viewed, similar to the original 0.1MHz signal sample.
Also, when the signal frequency is exactly equal to the sample frequency,
every data point taken will fall on the same position, forming a line.
\begin{figure}[H]
\begin{centering}
\includegraphics[scale=0.54]{/Users/ppkantorski/Documents/Radio_Astronomy/Lab_2/Pictures/power_leak}
\par\end{centering}

\caption{The graph above shows the frequency resolution and power leakage from
the 0.3MHz data sample used in Figures 1 and 2. The frequecy resolution
is half of the width of the spike.}
\end{figure}


Equation \ref{eq:freq_res} represents the frequency resolution of
a sampled waveform. Frequency resolution can be thought of as the
width between each frequency in the power spectrum. For the power
spectrum shown in Figure 4 and 5, the frequency resolution is approximately
$\Delta v\approx3.906kHz$.

\begin{equation}
\Delta v=\frac{1}{T_{sample}}=\frac{v_{sample}}{N}\label{eq:freq_res}
\end{equation}


Spectral leakage for the 0.3MHz data sample are shown in Figure 6.
The power spectrum is the average distribution of the power of a signal
represented as a set of components harmonically related to the fundamental
frequency, or the lowest frequency of the signal. When all the power
of the signal is at frequencies that are integer multiples of the
fundamental frequency, it can then adequately be represented by the
Fourier transformation. However, a problem arises when there is power
in the signal at frequency components that are not harmonics of the
fundamental. When this happens, the power of these components becomes
misrepresented. The term spectral leakage refers to this misrepresentation
of components that are not integer multiples of the fundamental frequency.


\subsection{Analog Mixing}

Using two SRS synthesizer oscillators, my partners and I DSB mixed
a 100kHz, -7dB$_{m}$ sine wave with a separate 105kHz and 95kHz,
-3dB$_{m}$ signal using a Mini-Circuits ZAD-1. Using the Pulsar card,
we chose sample at a rate of 4MHz and collect N = 16000 data points,
making $T_{sample}\approx4ms$. The resulting mixed signal is shown
in Figure 7 and power spectrum in Figure 8.

\begin{center}
\begin{tabular}{|c|c|c|c|}
\hline 
$v_{lo}$ & $v_{signal}$ & $v_{lo}\pm v_{signal}$ & $-v_{lo}\pm v_{signal}$\tabularnewline
\hline 
\hline 
100kHz & 105kHz & 205kHz, -5kHz & 5kHz, -205kHz\tabularnewline
\hline 
100kHz & 95kHz & 195kHz, 5kHz & -5kHz, 195kHz\tabularnewline
\hline 
\end{tabular}
\begin{figure}[H]
\begin{centering}
\includegraphics[scale=0.42]{/Users/ppkantorski/Documents/Radio_Astronomy/Lab_2/Pictures/analog_mix}
\par\end{centering}

\caption{This graph shows an analog DSB signal mixed between a 100kHz and 105kHz
signal obtained through Python with $v_{signal}=1.05v_{lo}$on the
left and $v_{signal}=0.95v_{lo}$on the right. A sample size of 16000
data points were sampled at a sample frequency of 4MHz using the Pulsar
card. There is also an apparent DC offset of about 0.01V. }
\end{figure}

\par\end{center}

As seen on Figure 7 and the table above, the mixed signal has four
distinct positive and negative frequencies. Additionally, since the
$\pm$5kHz signal is degenerate, there are only six unique frequencies.
For the positive $v_{lo}$ frequencies, the output splits into a lower
sideband$v_{lo}+v_{signal}$ and an upper sideband $v_{lo}-v_{signal}$,
forming the double sideband (or DSB) mixer. Single sideband (or SSB)
mixers require the mixing of pure sine waves that have both a real
and imaginary part. With complex sine waves, we are able to eliminate
one of the two sidebands. The upper and lower sideband of the mixed
signal are 205kHz and 5kHz signals respectively. The upper and lower
sideband are symmetric for the both negative and positive $v_{lo}$.
\begin{figure}[H]
\begin{centering}
\includegraphics[scale=0.42]{/Users/ppkantorski/Documents/Radio_Astronomy/Lab_2/Pictures/DSB_spec}
\par\end{centering}

\caption{The graph above shows the DSB mixing power leakage spectrum of the
mixed wave form in Figure 6. Dotted-line is placed along the distinctive
frequencies of the mixed signal. }
\end{figure}
\begin{figure}[H]
\begin{centering}
\includegraphics[scale=0.42]{/Users/ppkantorski/Documents/Radio_Astronomy/Lab_2/Pictures/filtered_analog}
\par\end{centering}

\caption{The graph above shows the inverse Fourier transform or Fourier filter
of the DSB mixing signal from Figure 6, isolating the 5kHz frequency.}
\end{figure}


Since the upper sideband distorts the 5kHz signal, I decided to filter
it out. To do this, using Python, I made the values of each frequency
higher on the FFT of the signal zero and then took the inverse FFT
(or IFFT) to retrieve filtered signal. This was done in Python by
first isolating 5kHz then using the ``IFFT'' feature built into
the package ``NumPy.'' This method of digital filtration is called
Fourier filtering and can be seen in Figure 9. The band of a Fourier
filter is somewhat arbitrary and really depends on some particular
range of interest. In the data, we chose also prevented any mixing
products that may have occurred as a result of any DC bias mixed into
our signals. After isolating the 5kHz and taking the inverse Fourier
transform of the power spectrum, the filtered signal was be retrieved
from the mix.


\subsection{Digital Mixing (DSB and SSB) on an FPGA}


\subsubsection{DSB Mixing}

Digital DSB mixing is similar to analog DSB mixing because of how
the signals are mixed. However, instead of mixing signals with physical,
analog components, the mixing is performed digitally in a computer.
But like analog mixing, there are still some inherit limitations.
Even if the two signals being mixed together do not violate the Nyquist
criteria individually, if the upper sideband of the signal is higher
than the Nyquist frequency, then aliasing will still occur. Furthermore,
if a DC offset is present in any of the separate signals, then it
will also show up as spikes on the power spectrum of the mixing product
around their frequencies. 

\begin{center}
\begin{tabular}{|c|c|c|c|}
\hline 
$v_{lo}$ & $v_{signal}$ & $v_{lo}\pm v_{signal}$ & $-v_{lo}\pm v_{signal}$\tabularnewline
\hline 
\hline 
1MHz & 1.05MHz & 2.05MHz, -50kHz & 50kHz, -2.05MHz\tabularnewline
\hline 
1MHz & 950kHz & 1.95MHz, 50kHz & -50kHz, -1.95MHz\tabularnewline
\hline 
\end{tabular}
\begin{figure}[H]
\begin{centering}
\includegraphics[scale=0.46]{/Users/ppkantorski/Documents/Radio_Astronomy/Lab_2/Pictures/digi_mix}
\par\end{centering}

\caption{The graph above shows a DSB mixed signal taken from the ROACH using
a LO frequency of about 100MHz and a signal frequency of about 95MHz.
The inverse Fourier transform of the signal is mapped on top of the
mixed signal, isolating 50kHz. Data was collected under a sample rate
of 200MHz.}
\end{figure}

\par\end{center}

As we have shown before, mixing of signals $v_{lo}$ and $v_{signal}$
is a fairly simple process. Digital mixing of two signals was performed
using the ROACH with our output results written to a BRAM file. Following
the same signal mixing procedures as Section 3.2 should produce similar
results when using the ROACH to digitally mix signals as shown in
Equation \ref{eq:analog_mix}. When data collection is triggered on
the ROACH, data points are collected and stored in BRAM until the
BRAM is full. When the BRAM is full, the number of data points taken
is constricted. BRAM files can only store 2048 ``32-bit signed''
integers and because the sample rate is 200MHz, our samples time was
about 10.24\textgreek{m}s. 
\begin{figure}[H]
\begin{centering}
\includegraphics[scale=0.42]{/Users/ppkantorski/Documents/Radio_Astronomy/Lab_2/Pictures/digital_DSB_power}
\par\end{centering}

\caption{The graph above shows the power spectrum of the digital of the DSB
mixing signal. The dotted lines show the 2.05MHz upper band on the
first graph, 50kHz lower band, and 1.95MHz band on the second graph.}
\end{figure}


Although we wanted to be consistency with out measurements, because
of problem with data collection, we used a local oscillator frequency
of 1MHz along with a 1.05MHz and 0.95MHz sine wave. Additionally,
the clock of the ROACH was set to at a rate of 200MHz. The clock sets
the rate at which logic gates in the FPGA execute operations and also
defines the rate at which data points are sampled. The results of
digitally mixing these signal frequencies can be seen in Figure 10.
The waveform and power spectrum are very similar to the results of
the analog mixing case, with the same spectral features. 

However, the signal appears wider than the the DSB result of analog
mixing using the Pulsar sampler (see Figure 7). The wider peaks are
caused by the limitations of our sampling time. When samples are taken
for longer periods of time, the frequency resolution of our spectrum
increases.The spectrum of the digital mixing (see Figure 11) is also
lot more smooth than the analog DSB mixing despite the DC offset.
This is because the digital mixing product only contains the mixed
frequency. In Figure 8, there are many noticeable higher-order frequency
components. Additionally, there appear to be less frequencies in the
spectrum of the digitally mixed DSB signal.


\subsubsection{SSB Mixing}

Although DSB mixing is useful, it can become problematic when mixing
involves a ``real'' local oscillator signal. Single sideband (or
SSB) mixing is one method of getting around this problem. Real signals
require a superposition of the sine wave incorporating both of the
positive and negative frequencies (see Figure 12). SSB mixing works
by taking a signal and splitting it into two identical signals, one
a cosine wave and another a sine wave. The cosine mixed signal is
then treated as the real component of the mixed signal while the mixed
sine wave represents the magnify components. The power spectrum of
this mix is shown in Figure 13. 

\begin{center}
\begin{tabular}{|c|c|c|}
\hline 
$v_{lo}$ & $v_{signal}$ & $v_{lo}\pm v_{signal}$\tabularnewline
\hline 
\hline 
625kHz & 1MHz & 1.625MHz, -375kHz\tabularnewline
\hline 
\end{tabular}
\begin{figure}[H]
\begin{centering}
\includegraphics[scale=0.42]{/Users/ppkantorski/Documents/Radio_Astronomy/Lab_2/Pictures/ssb}
\par\end{centering}

\caption{The graph above shows SSB digital mixing of a 10MHz complex signal
with the blue function representing the real part and the red representing
the imaginary. Data was collected under a sample rate of 200MHz.}
\end{figure}
\begin{figure}[H]
\begin{centering}
\includegraphics[scale=0.42]{/Users/ppkantorski/Documents/Radio_Astronomy/Lab_2/Pictures/ssb_power}
\par\end{centering}

\caption{The graph above shows SSB digital mixing of a 10MHz complex signal
with the blue function representing the real part and the red representing
the imaginary. Data was collected under a sample rate of 200MHz.}
\end{figure}

\par\end{center}

In Figure 12, the positive frequency (cosine wave) is indicated by
the real component of the filtered signal rising and falling slightly
out of phase with the imaginary component (sine wave) of the signal.
In both of the cosine and sine wave mixes, the frequency $v_{lo}-v_{signal}$
are the frequency of interest. Unlike DSB mixing, SSB mixing does
not mix signals in a superposition of frequencies. Because the power
spectrum is shifted with SSB mixing, concepts of ``upper'' and ``lower''
sidebands do not make sense. Since the input signal is real (see Equation
\ref{eq:complex}), there must be two distinctive frequencies in the
mixed product. Using a Fourier filter, the frequencies were isolated
from from the mixed signals. The sign of the lower frequency, however,
is determined from whether or not $v_{lo}$ is larger than the $v_{signal}$. 


\subsection{FIR Filter Response}


\subsubsection{FIR Coefficients}

The shape of the FIR filter chosen was determined by the desired bandpass
my partners and I specified. The impulse response of this bandpass
was set by the number of coefficients. The coefficients representing
the impulse response is called ``taps'' (see Figure 3). The order
of the filter we used was one less than the number of taps. As mentioned
earlier, the FIR filter we used for in the ROACH at our lab is an
8-tap FIR filter which we derived from writing the coefficient registers
to the FPGA.

\begin{table}[H] \center \begin{tabular}{|c|c|c|} \hline \multicolumn{3}{|c|}{Coefficients} \\ \hline Decimal Floating  &   Binary Fixed Point  & Hex representation \\ \hline $0.125$   & $0.00100000000000000$ & $\backslash x00\backslash x00\backslash x40\backslash x00$ \\ \hline $-0.0517767$ & $1.11110010101111101$ & $\backslash x00\backslash x03\backslash xe5\backslash x7d$ \\ \hline $-0.125$ & $1.11100000000000000$ & $\backslash x00\backslash x03\backslash xc0\backslash x00$ \\ \hline $0.3017767$ & $0.01001101010000010$ & $\backslash x00\backslash x00\backslash x9a\backslash x83$ \\ \hline $0.625$ & $0.10100000000000000$ & $\backslash x00\backslash x01\backslash x40\backslash x00$ \\ \hline $0.3017767$ & $0.01001101010000010$ & $\backslash x00\backslash x00\backslash x9a\backslash x83$ \\ \hline $-0.125$ & $1.11100000000000000$ & $\backslash x00\backslash x03\backslash xc0\backslash x00$ \\ \hline $-0.0517767$ & $1.11110010101111101$ & $\backslash x00\backslash x03\backslash xe5\backslash x7d$ \\ \hline \end{tabular} \caption{Coefficients used for the Digital Down Conversion on the   Roach. Each number after the backslash ($\backslash$) and $x$   represent the value of 4 bits, so each two numbers holds a byte. We   can see that summed they evaluate the 32 bits for the value of \emph{coefficient}.}  \label{tab:1} \end{table}

To choose coefficients for our FIR filter, we first designed our bandpass
in the frequency domain and then apply the IFFT of the bandpass to
retrieve our signal shape in the time domain. Since we did this digitally,
we chose 8 points on the filter so that the IFFT would generate a
complex coefficient for each of the eight registers. The filter we
constructed for this lab is called a 5/8-filter. We call this a 5/8-filter
because it was designed to filter the lower 5/8 portion of the frequencies
that the FFT can represent. Additionally, because filter is designed
to be symmetric about 0Hz, all of the filter's coefficients had to
be real. Since our FPGA's clock or sample rate is 200MHz, the Nyquist
frequency of our sampling is about 100 MHz. This made our desired
frequency range from about \textminus{}100MHz to 100MHz. The full
response of our 5/8-band filter is 5/8 of this range; from \textminus{}62.5
MHz to 62.5 MHz. The coefficients calculated for the 5/8 filter can
be found in Table 1.

The program we used on the ROACH was set to handle 18-bit fixed point
numbers. The numbers listed in the \textquotedblleft{}Decimal Floating
Point\textquotedblright{} column had to be converted to hex-digits
representing each bit for the fixed point numbers. This is done by
converting filter coefficients into an integer that have the same
bits as the 18-bit fixed point representation. Using the bits contained
in the integer, the hex-digits values were then used to control the
ROACH. Using Python, my partners and I generated both hex digits used
to control the ROACH along with the binary representations of the
filter coefficients. Because there are only 8 coefficients, the coefficients
alone do not provide a fine frequency resolution. To get the prediction
of the response of our filter have a ``finer'' frequency resolution,
my partners and I placed zeroes on the end of the list for the filter's
coefficients, virtually increasing our sample time.
\begin{figure}[H]
\begin{centering}
\includegraphics[scale=0.42]{/Users/ppkantorski/Documents/Radio_Astronomy/Lab_2/Pictures/FIR_filter}
\par\end{centering}

\caption{The graph above shows the FIR filter response using a -10 dBm, 100MHz
frequency noise source input. The local oscillated frequency was set
to 0Hz. The black line represents the filter's response. }
\end{figure}



\subsubsection{Testing FIR Filter Response}

After calculating our coefficients and producing a theoretical response
curve of the filter, the actual response of the filter can be examined
experimentally. By taking the noise 100MHz shown in Figure 14 and
applying the filter to it, the outline of the power spectrum can be
obtained. 
\begin{figure}[H]
\begin{centering}
\includegraphics[scale=0.5]{/Users/ppkantorski/Documents/Radio_Astronomy/Lab_2/Pictures/FIR_predic}
\par\end{centering}

\caption{The graph above shows the FIR filter response. Dotted line shows predicted
filter response, red points represent the coefficients, and solid
black line represents the response of the filter. }
\end{figure}


Although this is not really a precise measurement, but it provides
us with a rough estimate that the FIR filter we constructed is doing
what it is supposed to be. By taking the ratio of the filtered power
spectrum to the unfiltered power spectrum at its peak frequencies,
a more precise response of the FIR filter can be obtained. What Figure
15 shows is ratio between the square of the filtered to the square
of the unfiltered signal.


\section{Conclusion }

Signal processing is used in many applications and fields of study,
especially Radio Astronomy. Through testing the Nyquist criterion,
my partners and I found that frequencies up to one-half the sampling
rate can be accurently facsimliated. In this lab, we digitally sampled
signals, analog and digitally mixed signals, performed various methods
of filtration, and examined the FIR filter response set by the ROACH.

Additionally, my partners and I examined both analog and digital methods
of DSB and SSB signal mixing. We were also able to make comparisons
with the advantages and disadvantages of each. Through using techniques
like Fourier filtering, digital filtration was made possible. Since
digital mixing and signal filtration are used regularly in Radio Astronomy,
understanding just what exactly these signals are and how they can
be filtered is extremely relative.

In the end, I also learned about another method of digital filtering
in the down conversion process using a FIR filter. FIR filters allow
us to filter out unwanted frequencies from a DSB signal mixer. However,
we were limited by the number of taps or coefficients used to convolve
our signal. For purposes of isolating a certain frequency, the ``rippled''
shape of our the filter response (see Figure 15) will worked well
since the response eliminated frequencies outside our bandwidth. By
combining a FIR filter set with a specified sampling rate and using
a heterodyne mixer, radio signals can be converted into useful forms
for analysis.


\section{References}
\begin{enumerate}
\item Lab Partners: Leonardo Cassara and Caleb Levy.
\item Additional data was provided by Isaac Domagalski.
\item Help with data collection, analysis, and the report was also proved
by Baylee Bordwell, Eduardo Herrera, and especially Isaac Domagalski.
\item Undergraduate radio lab. (2014, March 7). Retrieved from


https://casper.berkeley.edu/astrobaki/index.php/Undergraduate\_Radio\_Lab.\end{enumerate}

\end{document}
